% Pista ana-24s-q5 · Anàlisi
% Procedència: PAU Catalunya 2024 · Convocatòria de setembre · Sèrie 3 · Qüestió 5
\documentclass[11pt,a4paper]{article}
\usepackage{pista}
\pistainfo{Catalunya 2024}{Convocatòria de setembre}{Sèrie 3}

\begin{document}

\section*{\textcolor{blauPista}{Qüestió 5}}

\noindent Per a cada punt $(x, y)$ de la corba $y = e^{-2x}$, amb $x > 0$ i $y > 0$, considereu el rectangle amb vèrtexs als punts $(0, 0)$, $(x, 0)$, $(0, y)$ i $(x, y)$.

\subsection*{\textit{a)} \normalfont Comproveu que, d'entre tots aquests rectangles, el que té $x = \dfrac{1}{2}$ és el d'àrea màxima. Quin és el valor d'aquesta àrea? \hspace*{\fill}\textnormal{[1,5~punts]}}

\pista{És un problema d'optimització.
\begin{itemize}
\item \textit{Funció a optimitzar.} El rectangle té base $x$ i alçada $y = e^{-2x}$. Escriu l'àrea en funció d'$x$: $A(x)$.
\item \textit{Optimitza.} Deriva $A(x)$ (recorda que és un producte de funcions) i resol $A'(x) = 0$. Comprova que $x = \dfrac{1}{2}$ és un màxim.
\end{itemize}

\medskip
\textbf{Atenció:} respon \emph{totes} les parts de la pregunta: també et demanen el valor de l'àrea màxima.}

\subsection*{\textit{b)} \normalfont Calculeu l'equació de la recta tangent a la funció $y = e^{-2x}$ en el punt d'abscissa $x = 0$, i el seu punt de tall amb l'eix de les abscisses. \hspace*{\fill}\textnormal{[1~punt]}}

\pista{Anomena $f(x) = e^{-2x}$. La recta tangent és
\[
y - f(0) = f'(0)\,x
\]
\begin{itemize}
\item Calcula $f'(x)$ (compte: no és la $A'(x)$ de l'apartat anterior) i avalua $f(0)$ i $f'(0)$.
\item Per al tall amb l'eix d'abscisses, imposa $y = 0$ a l'equació de la tangent i aïlla $x$.
\end{itemize}

\medskip
\textbf{Atenció:} la resposta final és un punt, amb les dues coordenades.}

\end{document}
